\textbf{Pour cette première partie, aucune justification n'est demandée et une seule réponse est possible par question. \\
Pour chaque question, reportez son numéro sur votre copie et indiquez votre réponse.}

Une réponse fausse ou l'absence de réponse n'enlève aucun point.

\subsubsection*{Question 1}

Le nombre $\left(5^{3}\right)^{2}$ est égal à :

\medskip

\noindent
\begin{tblr}{colspec = {X[c]X[c]X[c]X[c]},hlines,vlines}
\textbf{A.}\quad $5^5$
&
\textbf{B.}\quad $5^6$
&
\textbf{C.}\quad $5^9$
& \textbf{D.}
\quad  $5^1$
\\
\end{tblr}


\subsubsection*{Question 2}

Pour calculer le prix d'un produit après une hausse de 25~\%, en une seule opération, il faut multiplier le prix initial par :

\medskip

\noindent
\begin{tblr}{colspec = {X[c]X[c]X[c]X[c]},hlines,vlines}
\textbf{A.}\quad $0,25$
&
\textbf{B.}\quad $0,75$
&
\textbf{C.}\quad $1,25$
&
\textbf{D.}\quad  $1,75$
\\
\end{tblr}

\subsubsection*{Question 3}

Dans un repère, on a tracé la droite passant par l'origine du repère de coordonnées $(0 ; 0)$ et par le point de coordonnées $(2 ; 6)$.

\begin{center}
\begin{tikzpicture}
\begin{axis}[/pgf/number format/.cd, use comma,1000 sep={\,},
				width=12cm,height=10cm,
				xmin=-1.2,xmax=3.2,ymin=-3.2,ymax=9.2,
				xtick={-2,-1,...,4},
				minor x tick num =0,
				ytick={-4,-3,...,10},
				minor y tick num =0,
				tick label style={font=\small},
				grid=both, axis lines=center,
				grid style={gray!75,line width=0.4pt},
				minor grid style={gray!20,line width=0.2pt},
				enlargelimits=false]
	\addplot[color=blue,domain=-2:4,smooth,samples=2,line width=1.2pt] {(3*x)};
	\node[draw,cross out, line width=1.1pt,minimum width=3pt, minimum height = 3pt] at (axis cs:2,6){};
	\node[below right] at (axis cs:0,0){\small 0};
\end{axis}
\end{tikzpicture}
\end{center}

Une équation de la droite est :

\medskip

\noindent
\begin{tblr}{colspec = {X[c]X[c]X[c]X[c]},hlines,vlines}
\textbf{A.}\quad $y=2x+6$
&
\textbf{B.}\quad $y*6x+2$
&
\textbf{C.}\quad $y*3x$
&
\textbf{D.}\quad  $y=x$
\\
\end{tblr}


\subsubsection*{Question 4}

Dans un lycée, les adhérents de l'association sportive représentent un quart des élèves.

Quelle est la proportion d'élèves du lycée qui ne sont pas adhérents à l'association sportive ?

\medskip

\noindent
\begin{tblr}{colspec = {X[c]X[c]X[c]X[c]},hlines,vlines}
\textbf{A.}\quad 25~\%
&
\textbf{B.}\quad 75~\%
&
\textbf{C.}\quad 40~\%
&
\textbf{D.}\quad 60~\%
\\
\end{tblr}

\subsubsection*{Question 5}

Le nombre $\left(1+\dfrac{1}{2}\right)\times\dfrac{2}{3}$ est égal à :

\medskip

\noindent
\begin{tblr}{colspec = {X[c]X[c]X[c]X[c]},hlines,vlines}
\textbf{A.}\quad 1
&
\textbf{B.}\quad 2
&
\textbf{C.}\quad 3
&
\textbf{D.}\quad 4
\\
\end{tblr}



\subsubsection*{Question 6}

L'expression $(2x+3)^2$ est égale à :

\medskip

\noindent
\begin{tblr}{colspec = {X[c]X[c]X[c]X[c]},hlines,vlines}
\textbf{A.}\quad $2x^2+9$
&
\textbf{B.}\quad $4x^2+12x+9$
&
\textbf{C.}\quad $4x^2+9$
&
\textbf{D.}\quad $2x^2+12x+9$
\\
\end{tblr}


\subsubsection*{Question 7}

On lance un dé à 6 faces numérotées de 1 à 6.

On admet que toutes les faces ont la même probabilité d'apparaitre.

Quelle est la probabilité d'obtenir un numéro inférieur ou égal à 2 ?

\medskip

\noindent
\begin{tblr}{colspec = {X[c]X[c]X[c]X[c]},hlines,vlines}
\textbf{A.}\quad $\dfrac{1}{6}$
&
\textbf{B.}\quad $ \dfrac{2}{3} $
&
\textbf{C.}\quad $ \dfrac{1}{2} $
&
\textbf{D.}\quad $ \dfrac{1}{3} $
\\
\end{tblr}

\subsubsection*{Question 8}

On considère la série de nombres \quad $7 ; 3 ; 4 ; 2 ; x$.

La moyenne de cette série  de nombres est égale à 5.

Quelle équation permet de déterminer la valeur de $x$ ?

\medskip

\noindent
\begin{tblr}{colspec = {X[c]X[c]},hlines,vlines}
\textbf{A.}\quad $\dfrac{7+3+4+2+x}{5}=5$
&
\textbf{B.}\quad $7+3+4+2+\dfrac{x}{5}=5$
\\
\textbf{C.}\quad $\dfrac{7+3+2 + 4}{4}+x=5$
&
\textbf{D.}\quad $\dfrac{7+3+4+2+x}{4}=5$
\\
\end{tblr}
